
% 27 Mar Clean up 

\subsection{Installment 4: “Folkebladet” April 3, 1918}


% Sverdrup Between Lutheranism and Congregationalism [editorial title]


It will thus be seen that principle no. 1, "the proper chief principle of Scripture" (Veiledning, 143), is by no means a small matter. It is anything but a devout manner of speaking. Its rejection would, to a greater or lesser extent, modify the other fundamental propositions concerning church polity.

The Guiding Principles have again and again been attacked as being at variance with the New Testament’s teaching concerning the congregation. And in these attacks there has been no lack of proofs from Scripture. I mean that the thoughts which the Guiding Principles interpret, and the consequences which they involve, have fairly often and quite thoroughly been treated in the literature on church polity. The arguments have been both for and against—indirectly against, where Scripture has been broken in order to prove the supposed scriptural warrant for another concept of the church. For Norway’s part it is enough to point to a great number of articles in Luthersk Kirketidende, in Ny luthersk Kirketidende, and in Luthersk Ugeskrift. Germany has a large body of material to offer in its many works on ecclesiastical law and in innumerable contributions interpreting the most various church-political views. And the English Reformed literature on church polity is very extensive.

If we take up principle no. 1, we may first mention Professor von Oettingen (d. 1905) in Dorpat. He was known as a capable Lutheran dogmatician. But he unconditionally rejected the opinion that the New Testament knows only local congregations and contains no "doctrine of the church."

The teacher of ecclesiastical law, Professor Sohm, too, has several times, with good use of passages from Scripture, polemicized against that which principle no. 1 seeks to establish. As the best authority on Luther among jurists, he assures us that Luther could never have agreed with such a principle.

The church historian Professor Wm. Walter of Rostock, unquestionably the leading living authority on Luther among theologians, says the same of Luther and assures us that Jesus himself could not teach anything such as principle no. 1 teaches: it is too exclusive. (Ev. Luth. Kirchenzeitung, 1913.)

It is in the Reformed literature on church polity that principle no. 1 receives its strongest support. In the Lutheran it receives but little assent. I do not know of a single Lutheran theological faculty of rank in the whole world that would vote for it.

Nor do I know many Lutheran theologians (Professor Odland and Jens Greve must be counted among the exceptions) who agree that this principle is grounded in the New Testament. The great majority regard it as Reformed and will doubtless agree with the teacher of ecclesiastical law, Stahl, when he says: "The Reformed church has a strong tendency to let the concept of church pass over into the concept of congregation; it comes to that completely—in its offshoots, the Independents and Baptists, and so forth."

If one wishes to see how the principles of church polity in the Reformed church are judged by Lutherans, one does well to read Sohm’s Kirchenrecht, pp. 634–657, and especially the jurist Karl Rieter’s Grundsätze reformierter Kirchenverfassung (about 200 pages), a work that will cast a surprising light upon the Guiding Principles and their consistent interpretation.

But we need not stop with the German Lutherans. We may point to Norway. In 1915–16 a lively discussion concerning the concept of the church was carried on in Luthersk Kirketidende. It was between schoolmaster Jens Greve and parish pastor H. U. Sverdrup. Greve maintained the Free Church concept of the church, while H. U. Sverdrup rejected it. Both men are known as friends of the work for spirit and life.

Greve says:

\begin{quote}
"If we ask where the church is to be found on earth, then, according to our confession, we ask: Where is the Word preached purely, where are the Sacraments rightly administered? And where we find that this takes place, there we recognize the congregation to be. But then we turn our gaze to the local congregation; we cannot turn it elsewhere; even if it consisted of only two or three gathered in Jesus’ name, we must in it see the church made manifest.
\end{quote}

Greve teaches that the visible forms of manifestation are the church’s forms of manifestation, and that accordingly the true church of Jesus Christ in empirical reality comes forth in the visible congregations, there where the Word is preached purely and the Sacraments rightly administered. In other words: Not in an ecclesiastical organization embracing all Christians, nor yet in a mighty hierarchical order with a pope as supreme head, do we find the church, but in the local congregations, where the Word is rightly preached, etc., even where only two or three are assembled in Jesus’ name. We must acknowledge such an assembly, gathered around Word and Sacrament, as a congregation of Christ, a congregation of God; in such a congregation we see "the communion of saints," "the body of Christ," made manifest in a visible form of existence on earth."


This means: according to Scripture, the congregation is the right form of the kingdom of God on earth—which is precisely what the adherents of the state church and of the people’s church reject. Therefore Greve says:

\begin{quote}
"What, at bottom, is the fault among both the men of the state church and the adherents of the national church? They raise their building upon a foundation that fails; they build upon a false concept of the church," etc.
\end{quote}

To this H. U. Sverdrup replies:

\begin{quote}
"I maintain that in ordinary Norwegian this can only mean that in the constitutional question the men of the state church and the adherents of the national church are, broadly speaking, one, only that the men of the state church, or their theories, are perhaps a little worse in any case, for they are altogether impossible."
\end{quote}

Greve maintains that there is only one form of manifestation for the body of Christ. It is precisely the local congregation that is the body of Christ.

H. U. Sverdrup, on the other hand, maintains that there are many kinds of forms of manifestation. No visible form of existence can claim to be the body of Christ. He says:

\begin{quote}
"The Lutheran concept of the church has, within the evangelical division of the church, allowed itself to be joined with the most various forms of polity, the free-church forms as well, provided only that these, like this concept of the church itself, are of a folk-church character. The ‘free-church concept’ leads by inward consequence to the voluntaristic church with congregational order. For biblical and Lutheran thought, the legitimacy of the Free Church is dependent upon special conditions (the time of beginnings, the time of persecution, the end time); and since these conditions are always something that the church must seek to overcome, the Free Church (that is, the voluntaristic Free Church) will never stand as a goal that for principial reasons ought to be striven after. But this is precisely what the ‘free-churchmen,’ by virtue of their concept of the church, do. And indeed to such a degree that they designate the folk-church formation of the church as a deep fall."
\end{quote}

H. U. Sverdrup thus rejects that line of thought which makes the church’s outward organization something essential. Indeed he says:

\begin{quote}
"The concept of the free-churchmen is neither Lutheran, nor ancient-church, nor biblical, whereas the concept of the church among the adherents of the state church and of the folk church—precisely in the essential point that the church as the communion of saints, as the body of Christ, is invisible, and that it cannot in any visible form of existence claim to be this—agrees both with the Lutheran confession, the ancient-church conception, and with Scripture."
\end{quote}

Thus: the foregoing shows that there is fairly strong opposition to principle no. 1. It would be too lengthy to illuminate the other principles in this manner. Those who wish to have them illuminated will find them illuminated in the literature on church polity.

The Guiding Principles do not stand alone. But there lies a whole history behind them. They attempt to apply the Reformation principle of "the universal priesthood" to the local congregation and to the relation between the local congregations. The reason they are so little understood among us is doubtless the high-church upbringing which has entered into the blood even of many of the ablest and most zealous spokesmen of the revival: that direction which has set its stamp upon Krogh-Tonning’s doctrinal theology has also strongly influenced the home, religious instruction, the pulpit, the church press, indeed the whole of Christian thinking.

There are doubtless several directions that have been negatively determinative for the Guiding Principles. Thus the high-church tendency in Norway, with men like Heuch and Bugge at its head and with its many heirs in this country. We need not enter into these tendencies, but will only say that the the Guiding Principles have sought to take a stand against them.

It is of importance just now that we become more closely acquainted with the directions which, theoretically considered, have been positively determinative for the formulation of the principles and for the interpretation which their authors have given them.

---

If I am not mistaken, the authors have been influenced by four closely related currents:

1. The authors were strongly influenced by the church reform movement in Western Norway. In principle it recognized lay activity as something that normally belongs in an evangelical congregation. It maintained a democratic view of the ministerial office. It taught that all ecclesiastical right has its root in the congregation. In other words: "it championed the cause of religious liberty and congregational liberty" (Samlede Skrifter II, 201–211). This tendency found early expression in Luthersk Kirketidende. Later it received its organ in Ny luthersk Kirketidende, edited by Jakob Sverdrup and O. Vollan (1877–81). It was zealously opposed by Luthersk Ugeskrift, edited by Bugge and Heuch, who were then high-church.

To those who wish to know what polity is, we recommend Ny luthersk Kirketidende. Krogh-Tonning polemicized much against this paper and wrote of it in 1880: "The paper concerns itself almost exclusively with its favorite ideas about reforms in church polity; and it is precisely these that can only be explained by Reformed spirit, the spirit of independentism." He meant that the paper had fixed its gaze so blindly upon evangelical freedom that it had no sight for ecclesiastical authority; it had "a concept of the Christians’ spiritual priesthood which is at variance with the concept of the special priesthood"; it shut itself off from the truth that the office is instituted by God; it lacked the view of Jesus Christ’s church as the ‘one’ and ‘catholic’; and without further ado, in independentist fashion, transferred all right and power to the single local congregation. So much for Tonning.

2. Sverdrup and Oftedal nevertheless went further than the reform movement, for which the congregation was only the center of gravity in the ecclesiastical organization. They made the congregation the church itself, the only right form of the kingdom of God on earth. According to S. and O., the local congregation is an institution de jure divino.

But this is precisely the doctrine of the older Congregationalists (down to about 1870), although the newer school of Congregationalists maintains that there is no prescribed church form in the New Testament. To be sure, it finds no other church form in the New Testament than the local congregation, but it no longer constructs a de jure divino church polity on this basis. At the very most it regards the congregational form as something recommended, but not as something commanded. These are the ideas of the newer school, whereas S. and O. follow the old one with respect to the idea of the congregation’s divine institution.

The Congregationalists also went a step further than the reform party did in the democratic understanding of the ministerial office. They were very free with regard to lay activity and had no objection to a layman’s administering the Sacraments on the congregation’s behalf. But with respect both to the reception of members into the congregation and to church discipline, they were legal-minded, and therefore in this area diverged from S. and O., who here sought, as far as possible, to preserve evangelical freedom. In this respect the Congregationalists were more consistent than S. and O., who thought that logic must bend to life, although others object that principle no. 1 is a false premise, where life in reality bends to logic.

If S. and O. held the older Congregationalists’ view of principle no. 1 (the congregation is the right form) and of the equality of the members within the congregational body, they also shared their view with regard to the relation between the congregations. In his address "Free association among free congregations," 1897, Oftedal declared that the Lutheran Free Church, in the sense presented in his address, was "truly the old, good church order given of itself by the Word of God, in contrast to clerical rule, communal tyranny, and the state-church system."

How far S. and O. had acquainted themselves with the church polity of the Congregationalists (and Baptists) is not easy to determine. It is possible that they did not concern themselves much further with it, although Oftedal, as a church historian, must surely have known a good deal of it. It is possible that they took counsel with little besides the New Testament. In any case, Norway, despite the Haugean and Johnsonian revival, had not taught them much about the congregation. "There was no one in Norway," said Sverdrup, "who taught us anything about the free congregation or the true nature and form of Christ’s church, so that we all came over here rather unenlightened and without understanding in the things that belong to a free church fellowship" (S. S. II, 101).

But one thing must not be left unsaid. Oftedal and Sverdrup used the same arguments to prove the scripturalness of their conception of the congregation and of their view of the office, lay activity, and church organization as the Congregationalists did. The patterns of argument are remarkably similar, which they can only be when the presuppositions are the same. This applies not only to their marshalling of scriptural proof, but also to the polemic against the arguments which the opponents of the congregational principle have brought into the field. Everything converged on the claim that the congregation is Christ’s body, that it is therefore free and independent, that the other ecclesiastical organizations are human arrangements, and that the universal priesthood receives its application precisely in the congregation.
